%% kapa: exam=1 round=24 year=2013 period=1 form=A subject=civil_law subject_ko=민법 seq=13 number=13
%% source: Q-Net 감정평가사 기출문제 (cst00309) artlSeq=5073802 file=extracted/1008/5073802/2013년 제24회 감정평가사 1차시험 기출문제/1교시(A형).hwp
%% conversion: kordoc 4.24.51 -> tools/gen1.py (auto); review=proofread-sonnet; flags=none
\begin{question}{13}
복대리에 관한 설명으로 옳지 않은 것은?
\begin{choices}
  \item 법정대리인은 자신의 책임으로 복대리인을 선임할 수 있다.
  \item 복대리인은 임의대리인이다.
  \item 복대리인의 대리행위가 대리인의 대리권의 범위를 넘은 경우 본인이 이를 추인할 수 있다.
  \item 복대리인은 제3자에 대하여 대리인과 동일한 권리의무가 있다.
  \item 복대리인은 본인에게 직접 비용상환을 청구할 수 없다.
\end{choices}
\end{question}
\answer{5}
